\section{Mechanized Correctness in Lean 4}
\label{sec:correctness}

Compiler correctness is critical when applying aggressive global transformations. NumLang maintains machine-checked operational models and semantic preservation proofs in the \textbf{Lean 4} interactive theorem prover.

\subsection{Proof Architecture}
The formal proof codebase is organized across two Lake packages comprising 17 Lean 4 files totaling 1,699 lines:
\begin{itemize}
    \item \texttt{lean/Supercompiler/Semantics.lean}: Mid-Level IR (MIR) execution states, statement evaluation, and multi-step execution (\texttt{StepStar}).
    \item \texttt{lean/Supercompiler/Preservation.lean}: Multi-step driving transformation relations and evaluation preservation (\texttt{driving\_preserves\_semantics}).
    \item \texttt{lean/Supercompiler/Distillation.lean}: Global process-tree fold simulation and distillation relations (\texttt{distillation\_preserves\_semantics}).
    \item \texttt{lean/Supercompiler/MRSC.lean}: Pareto-optimal candidate selection preservation (\texttt{mrsc\_selection\_preserves\_semantics}).
    \item \texttt{lean/Supercompiler/Refinement.lean}: Path-sensitive interval pruning validity (\texttt{refinement\_pruning\_sound}).
    \item \texttt{lean/Supercompiler/Compaction.lean}: Control flow compaction and $\eta$-reduction (\texttt{noop\_removal\_preserves\_semantics}, \texttt{eta\_reduction\_preserves\_semantics}).
    \item \texttt{lean/Supercompiler/Main.lean}: Top-level composition theorem linking transformation passes (\texttt{supercompiler\_sound}).
    \item \texttt{proof/NumLangProofs/Driving.lean}: Core expression reduction, constant folding, branch pruning, and call unfolding (\texttt{smallstep\_const\_fold}, \texttt{smallstep\_call\_unfold}).
    \item \texttt{proof/NumLangProofs/Termination.lean}: Homeomorphic embedding preorder properties and finite-alphabet sequence pigeonhole termination (\texttt{finite\_configurations\_whistle\_terminates}).
\end{itemize}

\subsection{Main Correctness Theorem}
The central verification result in the MIR model establishes semantic preservation between the input program and the residual program:

\begin{theorem}[End-to-End Semantic Soundness (\texttt{supercompiler\_sound})]
\label{thm:soundness}
For any well-formed MIR program $p_1$ and transformed residual program $p_2$, if the supercompiler transformation sequence holds, evaluation outcomes are preserved:
\begin{equation}
\forall \text{args } \text{res}, \quad \text{Evaluates}(p_1, \text{args}, \text{res}) \leftrightarrow \text{Evaluates}(p_2, \text{args}, \text{res})
\end{equation}
\end{theorem}

\begin{proof}[Proof Structure]
The proof decomposes into inductive composition of verified phase relations:
\begin{enumerate}
    \item \textbf{Driving}: \texttt{driving\_preserves\_semantics} establishes that multi-step driving sequences preserve evaluation.
    \item \textbf{Folding \& Distillation}: \texttt{distillation\_preserves\_semantics} proves bidirectional simulation for knot-folding steps.
    \item \textbf{Compaction}: \texttt{noop\_removal\_preserves\_semantics} and \texttt{eta\_reduction\_preserves\_semantics} verify CFG compaction.
    \item \textbf{Path Refinement}: \texttt{refinement\_pruning\_sound} guarantees impossible interval branches cannot be taken.
\end{enumerate}
Composing these phase lemmas in \texttt{Main.lean} yields the top-level equivalence \texttt{supercompiler\_sound}.
\end{proof}

\subsection{Certified Integrity Guarantees}
Table~\ref{tab:lean_lemmas} summarizes the core mechanized theorems verified by the Lean 4 proof kernel.

\begin{table}[t]
\centering
\caption{Key machine-checked theorems in the NumLang Lean 4 formalization.}
\label{tab:lean_lemmas}
\small
\begin{tabular}{llr}
\toprule
\textbf{Module} & \textbf{Theorem Identifier} & \textbf{Formal Guarantee} \\
\midrule
\texttt{Semantics.lean} & \texttt{stepstar\_trans} & Transitivity of multi-step execution \\
\texttt{Preservation.lean} & \texttt{driving\_preserves\_semantics} & Driving step semantic preservation \\
\texttt{Distillation.lean} & \texttt{distillation\_preserves\_semantics} & Distillation fold simulation \\
\texttt{MRSC.lean} & \texttt{mrsc\_selection\_preserves\_semantics} & Candidate selection preservation \\
\texttt{Refinement.lean} & \texttt{refinement\_pruning\_sound} & Path constraint branch pruning \\
\texttt{Compaction.lean} & \texttt{noop\_removal\_preserves_semantics} & Dead block elimination soundness \\
\texttt{Driving.lean} & \texttt{smallstep\_const\_fold} & Small-step constant folding preservation \\
\texttt{Termination.lean} & \texttt{finite\_configurations\_whistle\_terminates} & Finite-alphabet whistle termination \\
\texttt{Main.lean} & \texttt{supercompiler\_sound} & \textbf{End-to-end semantic preservation} \\
\bottomrule
\end{tabular}
\end{table}

All proofs compile cleanly under Lean 4 (\texttt{lake build}) with zero custom axioms and zero \texttt{sorry} placeholders. As detailed in the documentation, these formal theorems verify abstract mathematical models of the compiler algorithms; the Rust executable itself is not extracted from Lean, and in-tree translation validation (\texttt{validate.rs}) is provided as an opt-in pass via the \texttt{--verify} flag.
